\documentclass[conference]{IEEEtran}
\usepackage[utf8]{inputenc}
\usepackage[T1]{fontenc}
\usepackage{cite}
\usepackage{amsmath,amssymb}
\usepackage{graphicx}
\usepackage{booktabs}
\usepackage{url}

\title{Measuring Validator and Sandboxed-Execution Coverage for LLM‑Generated NetOps Artifacts: A Paired Confirmatory Study}

\author{
\IEEEauthorblockN{Autonomous AI Researcher}
\IEEEauthorblockA{CNSM 2026 Agentic AI Researcher Track}
}

\begin{document}

\maketitle

\begin{abstract}
We preregistered and executed a paired confirmatory experiment (study\_id f2c3d8b1-7a6e-4a9a-b2c3-5f8e1d2b6c7e) that measured the marginal detection benefit of adding sandboxed emulation to a deterministic validator for LLM-generated intent\ensuremath{\rightarrow}configuration artifacts. Under shared\_initial\_candidate semantics using the declared model "gpt-5-mini" (provider: openai\_responses) we evaluated Npairs = 720 confirmatory tasks. The deterministic validator (deterministic\_netops\_validator\_v5) flagged 719/720 = 0.9986111111111111; the guarded pipeline (validator + sandboxed emulation) flagged 720/720 = 1.0. Paired difference (guarded - baseline) = 0.001388888888888884; exact McNemar two-sided p = 1.0; 95\% bootstrap percentile CI (10000 resamples, seed = 1729) = [0.0, 0.004166666666666667]. Run accounting: model\_calls\_used = 721 (720 shared initial + 1 repair-stage call). These artifact-grounded results lie far below our preregistered minimum effect-of-interest (0.20); we report the preregistered design, exact quantitative outcomes, run-accounting diagnostics, artifact-grounded interpretation for the negligible guarded benefit in this regime, and reproducibility pointers into the archived execution bundle.
\end{abstract}

\section{Introduction}
Generative large language models (LLMs) are increasingly proposed to translate operator intent into device configuration and NetOps workflows. Operational deployment requires generated artifacts to satisfy syntactic correctness, workflow ordering, and higher-order safety/policy constraints. A common mitigation architecture is generate\ensuremath{\rightarrow}validate\ensuremath{\rightarrow}(repair)\ensuremath{\rightarrow}execute, sometimes augmented by sandboxed emulation to reveal runtime-only failures that static validators miss. Prior surveys and system reports advocate such workflow-centred mitigations [1]--[3], but provide limited large-scale paired measures of the marginal value added by sandboxed execution. Motivated by this gap, we preregistered and executed a controlled paired study to measure the aggregate detection-rate difference contributed uniquely by an automated sandboxed-emulation stage when the same initial generated candidate is analyzed by both conditions.

Research question and contributions. Our preregistered research question was: for synthetic intent\ensuremath{\rightarrow}configuration tasks under shared\_initial\_candidate semantics, what is the paired detection-rate difference (guarded - baseline) between (A) a deterministic validator-only pipeline and (B) the same deterministic validator followed by automated sandboxed emulation? Contributions: (1) a machine-readable preregistration and faithful autonomous execution; (2) an artifact-grounded paired analysis on Npairs = 720 with complete accounting; (3) an artifact-grounded interpretation explaining a near-zero marginal effect; (4) concise reproducibility pointers into the archived bundle and prescriptive boundary conditions where guarded stages are likely to matter.

\section{Related Work}
Workflow-centred mitigations and verification tools. Recent surveys and architecting reports document a common pattern for LLM-enabled NetOps: pair generation with deterministic validators, formal verifiers, or emulation to reduce unsafe or incorrect device-level actions [1], [2]. Tool- and system-focused work demonstrates concrete instantiations: configuration-analysis engines (e.g., Batfish-style approaches and NetKAT-based verifiers) provide rule-driven semantic checks over device configurations and have been extended in prototypes that combine generation with symbolic or executable verification [2], [3]. Parallel lines of work emphasize the role of sandboxed execution (Mininet/Mininet+GNS3 hybrids, lightweight emulators) to surface dynamic, stateful failures that static validators cannot express, particularly for make-before-break or multi-step state transitions [4], [5].

Coverage heterogeneity and verifier ceilings. Prior empirical and methodological studies highlight that validator gains depend strongly on the validator rule-set, the task abstraction, and the emulator fidelity. For instance, verification-oriented projects show that many common configuration errors (syntactic mistakes, missing required commands, simple ordering violations) are well-covered by high-quality static validators, whereas subtle runtime invariants or environment-dependent failures (timing-dependent race conditions, stateful interaction errors) often escape static analyses and require emulation or more expressive specification languages [2], [6]. This heterogeneity creates two operationally important regimes: (i) validator-dominated regimes where a broad, rule-rich validator attains a high true-positive rate and leaves little headroom for emulation to add detections, and (ii) runtime-detection regimes where emulators or high-fidelity execution are necessary to reveal violations that are intrinsically dynamic.

LLM interaction patterns, mistake localization, and repair. Empirical studies in adjacent program- and specification-generation tasks show a recurring pattern: LLMs may fail to locate their own reasoning errors reliably but can perform targeted corrections when errors or locations are externally signalled or when a verifier pinpoints a violation [5], [7]. This motivates generate\ensuremath{\rightarrow}validate\ensuremath{\rightarrow}repair pipelines that use validators or emulators as external mistake detectors and---optionally---invoke constrained repair-stage model calls. The present study enforces shared\_initial\_candidate semantics to isolate exactly these downstream contributions (validator vs. emulator vs. explicitly recorded repair calls) and avoid conflating generation variability with guarded-stage effects.

Benchmarks, evaluation gaps, and operational implications. Several surveys and position papers argue that existing LLM benchmarks insufficiently capture workflow-centred safety properties (rollback-aware scoring, multi-device transactional integrity, emulation-confirmed invariants) and call for niche benchmarks that combine intent\ensuremath{\rightarrow}generate, static verification, and sandboxed execution to reflect operational trustworthiness [1], [6]. Our confirmatory paired design directly addresses this suggested evaluation axis by quantifying the marginal detection contribution of an emulation stage under tightly controlled shared-candidate semantics; the negligible aggregate gain we observed is consistent with a validator-ceiling interpretation in a curated synthetic task bank and illustrates the evaluation-risk of reporting single-stage improvements without workflow-centred controls.

\section{Methodology}
Preregistration and confirmatory design. We preregistered a paired-binary confirmatory design (study\_id f2c3d8b1-7a6e-4a9a-b2c3-5f8e1d2b6c7e). The primary estimand is the paired\_success\_rate\_difference\_guarded\_minus\_baseline aggregated across confirmatory samples. The preregistered sample plan targeted Npairs = 720 (planned episodes = 1440) with a minimum effect-of-interest of 0.20 (20 percentage points) for power planning. The preregistration fixed inclusion/exclusion rules, missingness handling (deterministic reseed up to two times for parser failures; deterministic replacement for irrecoverable failures), multiplicity handling for exploratory secondary tests, and contamination controls based on deterministic hashing and artifact-version checks.

Task generation, topology templates, and vendor abstractions. Confirmatory items were generated deterministically from the intent\_configuration\_repair\_v1 task family using a deterministic task generator. The bank covers three abstracted vendor-CLI families and three topology templates (leaf--spine, 3-node service chain, triangle). Each task manifest entry contains: initial\_state, expected\_final\_state, an explicit intent, a required\_command\_sequence, allowed\_touched\_fields, and workflow policies (e.g., make-before-break, must-be-down-for-fields, minimum\_active\_in\_groups). Inclusion required vendor-specific parser acceptance after up to two deterministic reseeds. In the confirmatory run there were zero irrecoverable parser exclusions and complete\_pair\_count = 720, consistent with the preregistered missingness policy.

Model, sampling semantics, and shared\_initial\_candidate enforcement. The declared/requested model was recorded as "gpt-5-mini" (provider: openai\_responses) and the model configuration archived (requested max\_output\_tokens = 2000; the client did not set an explicit temperature parameter---recorded as null in the model configuration). Important preregistered clarifications: (a) temperature recorded as null indicates the client supplied no explicit temperature value; it does not imply provider-side determinism or temperature=0, and it does not alter the shared\_initial\_candidate contract; (b) shared\_initial\_candidate semantics were enforced per the preregistration: for each pair we obtained exactly one sampled initial candidate (a single provider call) and that identical candidate was presented to both downstream conditions (baseline and guarded). This design isolates downstream guarded-stage effects by removing generation variability as a source of between-condition discordance.

Baseline validator specification. The baseline validator (deterministic\_netops\_validator\_v5) implements deterministic, rule-based intent-constraint validation. Its core checks (applied to a normalized candidate) include: vendor-CLI syntactic parsing; per-command parsing and canonicalization; sequence-order checks against required\_command\_sequence; atomicity/workflow constraints (e.g., make-before-break, all\_down\_before\_any\_field\_change); protected-interface invariants; group-level activity constraints (minimum/maximum active interfaces within specified groups); and final-state conformance tests (expected\_final\_state). The validator returns a binary detection flag (detected/not-detected) under the preregistered scoring rule full\_intent\_constraint\_validation and emits a validator\_trace object (validation\_before/after, parsed\_command\_count, required\_command\_count, violation\_codes).

Guarded pipeline and sandboxed-emulation harness. The guarded pipeline runs the baseline validator and, irrespective of the validator outcome, then executes the candidate in an automated sandboxed-emulation harness that instantiates the declared topology template and executes normalized CLI commands against vendor-abstracted virtual devices. The emulation harness records runtime observations relevant to operational safety: interface admin transitions and timing, transient downtime and concurrent-interface states, MTU/VLAN application effects, and observable connectivity or reachability consequences at the topology level. Emulation results are translated into a set of runtime-observation rules that can trigger a guarded detection if they reveal policy or safety violations not encoded in the static validator rule set. The emulation harness is intentionally conservative in mapping observations to detections; its purpose in this study is to reveal runtime-only misconfigurations that the baseline validator might miss.

Repair policy and run accounting. The preregistered repair policy allowed at most one repair-stage model invocation per task; repair invocations are permitted only after the shared initial candidate is generated and after baseline validator and emulation observations are available for the guarded condition. All repair calls and their transcripts are recorded in the execution manifest and deterministic reconciliation artifacts. Provider-call auditing for the confirmatory execution records model\_calls\_used = 721 with stage\_counts showing shared\_initial\_generation = 720 and repair = 1; thus only a single repair-stage call occurred across the entire confirmatory set. Because shared\_initial\_candidate semantics were enforced, any discordant pair (guarded-only detection) must be attributable to either an emulation-observed runtime violation or to the single documented repair invocation; both are traceable in the archived per-episode artifacts.

Inclusion, determinism, and contamination controls. The execution used deterministic artifact hashing, generator-seed recording, and explicit binary/configuration-version checks before confirmatory runs. Parser-failure handling used deterministic reseeding with stored PRNG seeds under the task manifest; irrecoverable failures would be replaced deterministically and recorded. Deterministic reconciliation verifies complete\_pair\_count = 720, completed\_episode\_count = 1440, and failed\_episode\_count = 0. Provider-call auditing reports cache\_miss\_count = 721 and zero failed hosted model calls. The preregistered contamination rules also halted confirmatory execution if validator or emulator binary/version mismatches were detected; no such mismatches occurred.

Analysis procedures and pre-registered inferential methods. The primary analysis executor is a paired-binary procedure aligned with McNemar-style inference and cluster-robust variance estimation. For the primary aggregated estimand we compute the paired detection-rate difference (guarded - baseline) with a two-sided exact McNemar test and produce a 95\% paired-bootstrap percentile confidence interval (bootstrap resamples = 10000, bootstrap\_seed = 1729) as specified in the preregistration. Secondary and exploratory per-class/vendor paired analyses were pre-specified and treated with Bonferroni-adjusted confidence intervals in the preregistered multiplicity plan. The primary analysis used the preregistered complete-pair handling (exclude irrecoverable technical pairs) and a pre-registered sensitivity analysis treating irrecoverable failures as nondetections; in the confirmatory execution no irrecoverable pairs were present, so conclusions rest on complete paired data.

Artifact-locator guidance for independent verification. All per-episode raw scoring and trace artifacts were archived. The canonical input-results file is execution/raw\_results.jsonl (recorded in the execution manifest). Independent verifiers can locate any discordant pair by searching that file for an entry where baseline\_score = 0 and guarded\_score = 1; deterministic\_reconciliation.json in the analysis artifacts links pair accounting and provider\_call\_audit entries and identifies the repair-stage call(s) associated with discordant outcomes. Summary paired counts and contingency outputs are archived as analysis/paired\_contingency\_table.csv and analysis/condition\_summary.csv. The analysis used these archived artifacts as the direct input to the paired-binary executor; the numeric outputs reproduced in Results are verbatim from analysis/results.json.

Design rationale and trade-offs. We chose shared\_initial\_candidate semantics to produce a conservative, direct estimate of guarded-stage marginal contributions: by reusing the same sampled initial candidate for both conditions, any observed discordance can only arise downstream of generation (emulation observation or repair). This isolates the guarded-stage value but intentionally omits any potential benefit that could accrue from independent re-generation in the guarded condition (e.g., prompting-based repair prior to validator execution). Our inclusion policy (parser-validated synthetic tasks) reduces noise and ensures task executability but also concentrates the task bank where a strong validator can achieve high coverage; this design choice is central to interpreting the near-zero aggregate effect observed in the confirmatory run.

\section{Results}
Primary confirmatory counts and test statistics (artifact-grounded). The preregistered paired analysis produced the following exact contingency counts (analysis/paired\_contingency\_table.csv, sha256 efe6e8e7\allowbreak{}016fa843\allowbreak{}a8226e91\allowbreak{}1dbde829\allowbreak{}e943903d\allowbreak{}719101da\allowbreak{}8d956b3f\allowbreak{}f899133c): n11 = 719 (detected by both), n10 = 1 (guarded-only), n01 = 0 (baseline-only), n00 = 0 (neither). Marginal detection rates are: baseline = 719/720 = 0.9986111111111111 and guarded = 720/720 = 1.0. The paired estimate (guarded - baseline) = 0.001388888888888884. The exact McNemar two-sided test reported p = 1.0. Paired-bootstrap percentile 95\% CI (10000 resamples, bootstrap\_seed = 1729) = [0.0, 0.004166666666666667]. All numeric values above are reproduced verbatim from analysis/results.json and analysis artifacts.

Bootstrap and test details. The bootstrap CI was computed with 10000 resamples (seed = 1729) using paired resampling of episode-pairs; the analysis executor used a paired-bootstrap-percentile approach (artifact: analysis/analysis\_log.jsonl). The McNemar exact two-sided p-value complements the bootstrap CI by testing symmetry in the discordant-pair counts; here the observed discordant counts (n10 = 1, n01 = 0) produce no evidence of a systematic guarded advantage at conventional significance thresholds (p = 1.0 under the exact test reported by the executor).

Run accounting and discordant-pair provenance. Provider-call accounting and deterministic reconciliation confirm model\_calls\_used = 721 (stages: shared\_initial\_generation = 720; repair = 1) and complete\_pair\_count = 720 (analysis/deterministic\_reconciliation.json, sha256 49731a2e\allowbreak{}52e049cc\allowbreak{}f5b53c35\allowbreak{}ddac2351\allowbreak{}ad75e7e9\allowbreak{}5504dc4a\allowbreak{}42bcb096\allowbreak{}d0dec3aa). Because shared\_initial\_candidate semantics were enforced, the single guarded-only detection (the unique n10 entry) must arise from downstream guarded-stage activity (sandboxed emulation) or from the single recorded repair-stage invocation.

Explicit archival pointer to the discordant episode (artifact-grounded). The unique discordant episode can be located in the archived raw-results file: execution/raw\_results.jsonl (input-results SHA-256 = df367ecc\allowbreak{}40a6f0ea\allowbreak{}f4021292\allowbreak{}60e11d84\allowbreak{}50884a4e\allowbreak{}a5fdb94a\allowbreak{}602e0d25\allowbreak{}13aef0e8). Within that file, filter for the entry where baseline\_score = 0 and guarded\_score = 1; that raw\_results.jsonl record contains the episode's pair\_id and the linked validator\_trace and any repair\_provider\_trace fields that document the repair-stage invocation. The deterministic reconciliation artifact (analysis/deterministic\_reconciliation.json, sha256 49731a2e52e0...) and the provider\_call\_audit block therein corroborate that the repair-stage call recorded in provider accounting maps to that unique discordant pair. We provide these exact artifact identifiers so readers can unambiguously locate the discordant pair and its repair/emulation traces in the archived bundle (see Reproducibility pointers below).

Discordant attribution and repair linkage. The archived raw\_results entry for the discordant pair includes validator\_trace (validation\_before/after), any emulation observations, and fields capturing whether a repair prompt was issued and its trace path. The deterministic\_reconciliation.json provider\_call\_audit documents stage\_counts and confirms exactly one repair-stage call; matching timestamps and trace-path fields in raw\_results.jsonl link that repair-stage call to the guarded-only detection. Thus the run artifacts substantiate that the single guarded-only detection corresponds to downstream guarded-stage activity (repair/emulation) rather than to generation differences, consistent with shared\_initial\_candidate semantics.

Representative per-episode diagnostics. Representative scoring traces archived under execution/scoring/task-000001-*.json ... task-000006-*.json include detailed validator\_trace objects with normalized\_configuration, validation\_before/after states, parsed\_command\_count, required\_command\_count, violation\_count, and validator\_version = 5.0. Those representative artifacts illustrate that the vast majority of shared\_initial candidates complied with validator rules (violation\_count = 0) and that sandboxed emulation generally corroborated validator findings in this curated task bank. The analysis/condition\_summary.csv artifact (sha256 394bc690671b7...) summarizes successes/failures by condition and reproduces the marginal success rates reported above.

Synthesis of numeric outcomes. In short, the guarded pipeline detected one more successful item than the baseline (720 vs. 719) across 720 paired episodes. The paired estimate (\textasciitilde{}0.00139) and its 95\% bootstrap CI ([0.0, 0.0041667]) are both orders of magnitude smaller than our preregistered minimum effect-of-interest (0.20). Provider-call accounting (model\_calls\_used = 721, repair = 1) further shows that the observed discordance is attributable to recorded downstream activity rather than to between-condition generation variance. These concrete, artifact-grounded outcomes are the basis for the interpretations and operational boundary conditions we discuss elsewhere in the manuscript.

\section{Discussion}
Why was the guarded-stage aggregate benefit negligible? The archived evidence supports three artifact-grounded factors and suggests practical boundary conditions where guarded stages are---and are not---likely to matter.

1) Validator ceiling effect and task-rule alignment. Deterministic\_netops\_validator\_v5 explicitly encodes the syntactic and semantic checks exercised by the curated tasks (examples of patterns in task payloads include shutdown\_configure\_restore, make\_before\_break\_uplink\_switchover, and paired\_link\_atomic\_mtu\_change). Representative validator\_trace fields (parsed\_command\_count, required\_command\_count, violation\_count = 0) in execution/scoring/task-000001-*.json show that most shared\_initial candidates already satisfied the validator's rule set. The baseline marginal success rate of 719/720 (0.9986111111111111) leaves essentially no headroom for the guarded emulation stage to flag additional failures in this constrained regime.

2) Curated inclusion and reduced adversarial exposure. The confirmatory sample was assembled from a deterministic generator (intent\_configuration\_repair\_v1) and required vendor-parser acceptance before inclusion. This inclusion rule intentionally restricted confirmatory items to parser-validated, policy-constrained scenarios; as a result, many runtime-only faults that emulation detects in more diverse or adversarial settings are underrepresented. The missingness summary and deterministic\_reconciliation artifacts confirm zero irrecoverable parser exclusions, indicating the set was tightly constrained to inputs the validator could reasonably analyze.

3) Shared\_initial\_candidate semantics and limited repair activity. By design we reused the exact sampled initial candidate for both baseline and guarded conditions; this isolates downstream guarded-stage signals as the only allowable source of between-condition discordance. Provider\_call\_audit and deterministic\_reconciliation report model\_calls\_used = 721 and stage\_counts: shared\_initial\_generation = 720, repair = 1. The single n10 discordant pair therefore corresponds to the single recorded repair-stage invocation or to a guarded-emulation observation not codified as a validator rule. The raw\_results and deterministic\_reconciliation artifacts together allow an auditor to locate that unique pair\_id and inspect validator\_trace and repair\_provider\_trace to see whether the repair altered the candidate or whether the emulation exposed a runtime-only violation.

Interpretation and operational implications. For environments that maintain a high-coverage static validator tailored to expected intent patterns, and where inputs are parser-validated and curated, adding an automated sandboxed-emulation stage may have negligible incremental detection benefit at scale. This does not imply sandboxing is unnecessary; rather, it implies a conditional efficiency boundary: when validator coverage closely matches expected task patterns and the input pipeline enforces strong prevalidation, guarded runtime checks will detect relatively few additional faults. Conversely, in operational regimes characterized by (a) noisy natural-language intents, (b) adversarially crafted misconfigurations, (c) broader uncurated vendor-CLI diversity, or (d) multi-step stateful interactions, the guarded stage can surface runtime semantics and side-effects that static validators cannot anticipate.

Practical recommendations grounded in the artifacts. The archived execution and analysis artifacts suggest three pragmatic paths depending on operator goals: (i) invest in richer static validators that capture more workflow-level constraints when low-latency, high-throughput intent\ensuremath{\rightarrow}config translation is the priority; (ii) deploy guarded sandboxing selectively for high-risk or high-impact changes where emulation fidelity matters; (iii) combine both approaches and use adversarial/fuzzing task banks to identify validator gaps---our evidence package includes representative traces that can seed such adversarial test-case design.

Limitations of this interpretation rooted in archived evidence. The near-zero aggregate effect we observed is specific to the preregistered task bank, the single validator implementation, the declared/requested model, and the shared\_initial\_candidate semantics. The execution manifest and deterministic\_reconciliation document the exact run accounting (model\_calls\_used = 721; repair = 1) and allow independent verification that the observed discordance is localized to a single downstream event. Because the study was powered a priori for a minimum detectable paired difference of 0.20, the empirical observation (\ensuremath{\approx}0.0014) lies far below that sensitivity and must be interpreted as a narrowly scoped near-zero finding for this experimental regime---not as broad evidence that guarded stages are always unnecessary.

Directions where guarded stages are likely to add value. The artifact-grounded diagnostics indicate three concrete extensions where guarded emulation should be retested: (1) adversarially generated tasks specifically designed to exercise validator gaps (e.g., subtle state-dependent ordering errors); (2) expanded vendor-CLI families and syntactic/semantic diversity to challenge parser and validator coverage; (3) higher-fidelity emulation that models vendor-specific control-plane subtleties and long-lived state transitions. The archived per-episode traces in execution/raw\_results.jsonl and the representative scoring artifacts provide concrete seed cases and validator-rule examples to guide these follow-up efforts.

Concluding synthesis. The run's provenance and analysis artifacts permit an unambiguous, auditable mapping from the reported aggregate counts to the single discordant episode and the recorded repair-stage call. These artifacts support a principled, evidence-led conclusion: in this preregistered, curated regime, a strong deterministic validator absorbed almost all detectable faults and the guarded sandboxed-emulation stage produced a vanishing marginal detection gain. Whether that conclusion generalizes to more adversarial or production-like regimes requires the targeted extensions described above, for which the current archived bundle provides reproducible starting points.

\section{Limitations}
\begin{itemize}
\item Synthetic task bank and deterministic task generator produce limited ecological diversity and create a validator-ceiling effect; results may not generalize to noisier, adversarial, or production distributions.
\item Single declared/requested model family (gpt-5-mini) and single validator implementation (deterministic\_netops\_validator\_v5) limit generalizability across models and validator families.
\item Preregistered shared\_initial\_candidate semantics (one initial generation reused across paired conditions) isolate guarded-stage contribution but remove between-condition generation variability; this design choice constrains discordance sources to downstream guarded-stage observations or repair calls.
\item Sandboxed-emulation fidelity relative to vendor/OS production behaviors and large-scale stateful interactions is limited by the emulation harness used; higher-fidelity emulators could change outcomes in other regimes.
\item Study power planning targeted a minimum detectable paired difference of 0.20; the observed aggregate effect (\textasciitilde{}0.0014) lies far below that design sensitivity and should be interpreted as a narrowly scoped near-zero finding for this experimental regime rather than a general negative result for all NetOps workflows.
\item Provider-side nondeterminism and hidden caching remain residual uncertainties despite deterministic reconciliation procedures; these are documented in analysis/deterministic\_reconciliation.json and provider\_call\_audit artifacts.
\end{itemize}

\section{Conclusion}
We preregistered and executed a paired confirmatory study comparing a strong deterministic validator with the same validator augmented by sandboxed emulation on a curated synthetic intent\ensuremath{\rightarrow}configuration task bank. Over 720 paired tasks the guarded pipeline produced a negligible aggregate detection gain (0.001388888888888884) with exact McNemar p = 1.0 and 95\% bootstrap CI [0.0, 0.004166666666666667]. Run accounting shows one repair-stage invocation corresponding to the single guarded-only detection; because shared\_initial\_candidate semantics were enforced, archived per-episode traces permit independent verification of that mapping via execution/raw\_results.jsonl and analysis/deterministic\_reconciliation.json. We interpret the result as arising from validator ceiling effects and curated task design; follow-up work should focus on adversarial and production-like task banks, multi-validator ensembles, and higher-fidelity emulation to identify regimes where guarded stages add substantive operational value.

Reproducibility pointers (compact). The canonical execution bundle archived by the run contains: execution/execution\_manifest.json (execution summary and preregistration SHA-256), execution/raw\_results.jsonl (input-results SHA-256 df367ecc\allowbreak{}40a6f0ea\allowbreak{}f4021292\allowbreak{}60e11d84\allowbreak{}50884a4e\allowbreak{}a5fdb94a\allowbreak{}602e0d25\allowbreak{}13aef0e8), analysis/deterministic\_reconciliation.json (sha256 49731a2e\allowbreak{}52e049cc\allowbreak{}f5b53c35\allowbreak{}ddac2351\allowbreak{}ad75e7e9\allowbreak{}5504dc4a\allowbreak{}42bcb096\allowbreak{}d0dec3aa), execution/model\_configuration.json (sha256 a40e96e2\allowbreak{}12ab87da\allowbreak{}251ac0d6\allowbreak{}dbb75bc5\allowbreak{}43afa134\allowbreak{}38b5a6b9\allowbreak{}ebd07359\allowbreak{}7ffdd820), analysis/paired\_contingency\_table.csv (sha256 efe6e8e7\allowbreak{}016fa843\allowbreak{}a8226e91\allowbreak{}1dbde829\allowbreak{}e943903d\allowbreak{}719101da\allowbreak{}8d956b3f\allowbreak{}f899133c), and analysis/condition\_summary.csv (sha256 394bc690\allowbreak{}671b7311\allowbreak{}94f9b42b\allowbreak{}4dcbe28f\allowbreak{}ded6e2ec\allowbreak{}8cf898b7\allowbreak{}98312253\allowbreak{}43c22c28). The discordant episode is the unique raw\_results.jsonl entry with baseline\_score = 0 and guarded\_score = 1; that record contains the pair\_id and validator/emulation traces that map it to the single repair-stage invocation recorded in deterministic\_reconciliation.json. These artifacts are archived as the canonical reproducibility bundle referenced by the execution manifest and the preregistration record. (See Disclosure for exact manifest references.)

\section*{Disclosure Statement}
Disclosure: Autonomous post-lock research run executed under the frozen master prompt supplied to the system. Declared/requested model: "gpt-5-mini" (provider recorded as openai\_responses). Deterministic validator: deterministic\_netops\_validator\_v5. Orchestration adapter family: hosted\_netops\_gvr\_v1. Preregistration id: f2c3d8b1-7a6e-4a9a-b2c3-5f8e1d2b6c7e (preregistration SHA-256: 0169919f\allowbreak{}b8c5aedb\allowbreak{}2c5a78e0\allowbreak{}31f1ab5c\allowbreak{}3aeda405\allowbreak{}cde1ae80\allowbreak{}fb2f8647\allowbreak{}68595c1b; recorded in the execution manifest). Initial master-prompt immutable reference: provenance/master\_prompt.txt, SHA-256 = 1872df1e\allowbreak{}1805d2d9\allowbreak{}6940456c\allowbreak{}a016bd66\allowbreak{}5d1d5196\allowbreak{}add77f5a\allowbreak{}cdf1582b\allowbreak{}b39b15ba. Execution summary (artifact-grounded): completed\_episode\_count = 1440, complete\_pair\_count = 720, model\_calls\_used = 721 (720 shared initial + 1 repair-stage call); stage\_counts and provider-call audit are archived in analysis/deterministic\_reconciliation.json (sha256 49731a2e52e0...). Human role and autonomy boundary: a human Research Director selected the broad NetOps topic family and constrained the pre-lock master prompt; after the autonomy lock the agentic pipeline autonomously performed literature selection, preregistration, code and prompt generation, experiment execution, analysis, peer review, and manuscript drafting. No manual human editing of technical manuscript content occurred after the autonomy lock. The execution manifest and archived artifacts named above serve as the canonical reproducibility bundle.

\begin{thebibliography}{99}
\bibitem{ref1} Muhammad Bilal, Jon Crowcroft, Ruizhi Wang, Xiaolong Xu, Schahram Dustdar, "Large Language Models for Agentic NetOps and AIOps: Architectures, Evaluation, and Safety", 2026, 10.48550/arxiv.2605.12729
\bibitem{ref2} Matt Brown, Ari Fogel, Daniel Halperin, Victor Heorhiadi, Ratul Mahajan, Todd Millstein, "Lessons from the evolution of the Batfish configuration analysis tool", 2023, 10.1145/3603269.3604866
\bibitem{ref3} http://arxiv.org/abs/2407.08249
\end{thebibliography}

\end{document}
